%======================================================================
% Section 6: Generalized Transport and Memory Effects
%======================================================================
\section{Generalized Transport and Memory Effects}
\label{sec:transport}

Classical Fourier conduction relates heat flux instantaneously to the temperature gradient. Its range of adequacy and any departure from it depend on material, scale, and experiment; no universal failure threshold is asserted here. This section surveys alternative constitutive structures and identifies questions that require model-specific derivation or source checking.

\subsection{Physical origin of thermal relaxation}

Microscopic heat carriers and their scattering can motivate finite-rate transport models, but a continuum relaxation parameter is phenomenological unless it is derived or identified from a specified kinetic model. An exponential relaxation law is one possible constitutive choice, not a universal consequence of carrier dynamics. The memory kernel, relaxation parameters, and resulting wave/support properties must be established for the particular model and material. This section surveys such operator choices while marking family-wide physical interpretations as conditional or unverified.

\subsection{Three structural routes to memory}

Three mathematically distinct ways of making a constitutive law remember have been used, and they are not interchangeable.

\begin{itemize}
\item \textbf{Differential-rate constitutive laws.} Some models include time derivatives of fluxes, temperatures, or other fields; the order and choice of differentiated variable are model-specific. Whether such a law admits an equivalent finite-dimensional internal-variable representation must be established for the particular formulation.
\item \textbf{Integral (convolution) memory.} The constitutive relation is written as a convolution of the driving quantity with a memory kernel over a finite or semi-infinite interval. This includes the hereditary-integral formulation of Section~\ref{sec:transport} and the memory-dependent derivative of Wang and Li \cite{Wang2011}. The memory is then encapsulated in an arbitrary function, and the price is that the constitutive relation is no longer local in time, so that the discretisation must carry the history.
\item \textbf{Spatial nonlocality.} Some constitutive laws average over a material neighbourhood, as in Eringen-type formulations. This is a spatial interaction, not temporal memory, and its operators and boundary conditions require separate treatment in Section~\ref{sec:nonlocal}.
\end{itemize}

\subsection{A transform-domain example, not a universal representation}

Some linear flux laws can be written schematically as
\begin{equation}
\mathcal{D}[\bm{q}] = -k\,\bm{\nabla}T,
\label{eq:6.1}
\end{equation}
but this notation does not make distinct constitutive models equivalent. For the Lord--Shulman/Cattaneo law, with zero initial heat flux,
\begin{equation*}
\bm q+\tau_0\dot{\bm q}=-k\nabla T.
\end{equation*}
The corresponding Laplace-domain multiplier is
\begin{equation}
\mathcal{D}_{\rm LS}(s)=1+\tau_0s.
\label{eq:6.2}
\end{equation}
Nonzero initial flux contributes an additional initial-data term. Fractional, phase-lag, and hereditary models have their own transform symbols and initial terms, which must be derived from their exact constitutive equations; no universal symbol or shared limiting procedure is asserted here.

A comparison in the transform domain is useful only after each constitutive law, initial condition, and admissible parameter range is stated. Formal limits such as $\tau_0\to0$, fractional order tending to an integer, or memory horizon tending to zero require model-specific normalization and convergence arguments. Positivity of a frequency-domain loss modulus, when defined, is not by itself proof of a consistent free energy or a well-posed initial-boundary-value problem. These questions remain model-specific and are not settled uniformly in this review.

\subsection{Fractional operators: Caputo and Riemann--Liouville}
\label{sec:6.2}

Fractional derivatives generalise integer-order differentiation to non-integer orders and are used in models with history dependence. Depending on the operator, coupling, kernel, and initial-boundary conditions, they can produce subdiffusive, dispersive, or wave-like responses. A finite propagation speed or a specific support property is not implied by the fractional order alone.

\subsubsection{Caputo fractional derivative}

The Caputo derivative of order $\alpha$ ($n-1 < \alpha < n$, $n = \lceil\alpha\rceil$):

\begin{equation}
{}^C D_t^\alpha f(t) = \frac{1}{\Gamma(n-\alpha)}\int_0^t \frac{f^{(n)}(\tau)}{(t-\tau)^{\alpha+1-n}}\,d\tau.
\label{eq:6.3}
\end{equation}

\subsubsection{Riemann--Liouville fractional derivative}

\begin{equation}
{}^{RL} D_t^\alpha f(t) = \frac{1}{\Gamma(n-\alpha)}\frac{d^n}{dt^n}\int_0^t \frac{f(\tau)}{(t-\tau)^{\alpha+1-n}}\,d\tau.
\label{eq:6.4}
\end{equation}

For suitable orders and regularity, the Caputo definition uses integer-order initial data, while Riemann--Liouville formulations may involve fractional initial data. This difference can matter in model selection, but it does not establish a universal preference. The historical motivations and material interpretations attributed to fractional operators require claim-level source checking; the citations here are source leads, not independently verified support for those claims.

\subsubsection{Fractional heat conduction and its limits}

\begin{equation}
{}^C D_t^\alpha T = a_\alpha\nabla^2 T, \qquad 0 < \alpha \leq 2,
\label{eq:6.5}
\end{equation}
Here $a_\alpha$ is a model coefficient with dimensions consistent with the selected order (for length and time units, $[a_\alpha]=L^2T^{-\alpha}$). This scalar equation is illustrative, not a universal fractional thermoelastic law. At the integer-order endpoints, the corresponding scalar PDE has heat-equation or wave-equation form, with coefficients and initial data specified for each model; one cannot infer a common dimensional parameter or solution limit across orders. Intermediate-order propagation and well-posedness depend on the operator definition, coupling, domain, and data. The cited Povstenko papers are relevant source leads for fractional thermoelasticity, but no theorem or material interpretation from them has been checked in full text in this revision.

\subsection{Hereditary integrals: convolution kernels}
\label{sec:6.3}

One common representation of linear temporal memory is a hereditary integral:

\begin{equation}
f(t) = \int_{-\infty}^{t} K(t-\xi)\,g(\xi)\,d\xi,
\label{eq:6.6}
\end{equation}
where $K(t-\xi)$ is a kernel in this schematic convolution. Some models impose fading-memory, positivity, monotonicity, or complete-monotonicity hypotheses, but the required conditions depend on the constitutive law and the chosen thermodynamic framework; none of these properties alone is asserted here as a universal admissibility criterion. For suitable completely monotone kernels, representation results by weighted exponentials may apply, but their hypotheses and implications must be checked for the kernel at hand. No equivalence with a discrete phase-lag hierarchy is claimed.

\subsection{Memory-dependent derivative (MDD): operator and kernel}
\label{sec:6.4}

An operator attributed to Wang and Li \cite{Wang2011} uses a kernel over a finite interval. The exact definition and its relation to other memory models remain subject to primary-source checking:

\begin{equation}
D_\omega f(t) = \frac{1}{\omega}\int_{t-\omega}^{t} \kappa(t-\xi) f'(\xi)\,d\xi,
\label{eq:6.7}
\end{equation}
where $\omega>0$ denotes the displayed horizon and $\kappa$ is the kernel. The bounds, normalization, regularity, and admissibility conditions on $\kappa$ must be taken from the exact model definition. A finite integration horizon truncates the history in this formula; it does not by itself imply a finite-dimensional state-space realization, finite computation cost, or that the kernel is a freely chosen rather than material-calibrated quantity.

\subsubsection{MDD kernel function}

A commonly used kernel is
\begin{equation}
\kappa(t-\xi) = 1 - \frac{2\beta}{\omega}(t-\xi) + \frac{\alpha^2}{\omega^2}(t-\xi)^2,
\label{eq:6.8}
\end{equation}
where $\alpha$ and $\beta$ are parameters in this inherited candidate kernel. Its dimensions, bounds, normalization, and any admissible parameter region require checking against the defining source and the selected thermodynamic model; positivity or monotonicity alone is not taken as a proof of admissibility.

\subsubsection{Limiting behavior}

A formal short-horizon limit $\omega\to0$ is often used to recover an ordinary derivative from a normalized finite-horizon operator. The displayed kernel normalization and regularity assumptions must be specified before this limit is asserted; it is retained as a candidate limiting relation and is not used as a proven model equivalence. No exact exponential-kernel reduction to Lord--Shulman, or general equivalence with a fractional operator, is asserted here.

\subsection{Energy flux in memory-dependent media}
\label{sec:6.5}

A candidate outward-positive local flux for a classical mechanical balance with a specified heat-flux vector is
\begin{equation}
J_i(t) = q_i(t)-\sigma_{ij}(t)v_j(t).
\label{eq:6.10}
\end{equation}
This is not asserted as a complete total-energy flux for every memory-dependent theory. The MDD heat-flux law and the appropriate stored-energy/boundary-power terms must be derived together from the selected kernel and free energy. The subsequent equations are inherited candidate constitutive forms, not independently source-checked here:
\begin{equation}
q_i(t) + \tau_0 D_\omega q_i(t) = -k T_{,i}(t),
\label{eq:6.11}
\end{equation}
\begin{equation}
q_i(t) + \frac{\tau_0}{\omega}\int_{t-\omega}^{t} \kappa(t-\xi) \dot{q}_i(\xi)\,d\xi = -k T_{,i}(t).
\label{eq:6.12}
\end{equation}
No universal hereditary correction to the energy flux, or generalized local balance, is asserted. Status of Eqs.~\eqref{eq:6.11}--\eqref{eq:6.12} and the corresponding MDD free-energy balance: UNVERIFIED / OPEN pending primary-source and derivation audit.

\subsection{Well-posedness and propagation in memory-dependent models}
\label{sec:6.6}

Well-posedness and propagation depend on the exact memory operator, kernel, time horizon, initial data, domain, and boundary conditions. Positivity or monotonicity of a kernel alone is not a general sufficient proof of a well-posed initial-boundary-value problem. Likewise, a finite-horizon derivative does not by itself imply finite propagation speed. The inherited claims about universal admissibility restrictions, a Banach-space existence theorem, and finite thermal-wave speed have not been independently verified and are not asserted as results in this revision. A model-specific energy estimate and existence/uniqueness analysis are OPEN.

\subsection{Comparative summary: transport models}

Table \ref{tab:transport_comparison} summarizes transport and memory models across mathematical structure, dissipation, well-posedness, and energy flux status.

\begin{table}[htbp]
\centering
\small
\caption{Provisional transport-family comparison. Entries for generalized models are audit questions, not verified admissibility or well-posedness results.}
\label{tab:transport_comparison}
\begin{tabularx}{\textwidth}{@{}p{2.4cm}XXX@{}}
\toprule
\textbf{Model class} & \textbf{Transport structure} & \textbf{Energy/admissibility status here} & \textbf{Evidence status} \\
\midrule
Classical Fourier & Local Fourier heat flux & Classical local balance derived; Fourier heat equation parabolic & DERIVED / standard model \\
Lord--Shulman & Relaxing heat flux & Model-specific free energy, dissipation and initial data not checked & OPEN \\
Green--Lindsay & Temperature-rate-dependent formulation & Single-temperature description; source-level equations and admissibility not checked & UNVERIFIED \\
Phase-lag (dual/triple) & Lagged constitutive laws & Restrictions and stability depend on selected truncation and parameter set & OPEN \\
Green--Naghdi II/III & Thermal-displacement formulations & Family-specific energy and dissipation identities not checked & OPEN \\
Moore--Gibson--Thompson & Higher-order heat equation variants & Stability and admissibility not checked here & OPEN \\
Memory-dependent & Kernel/time-history operator & Kernel-specific balance and well-posedness not checked & OPEN \\
Fractional & Fractional time operator & Propagation and admissibility depend on operator and data space & OPEN \\
Nonlocal / gradient & Spatially nonlocal or higher-gradient law & Boundary power and domain-specific balance not derived & OPEN \\
\bottomrule
\end{tabularx}
\end{table}

\begin{figure}[htbp]
\centering
\includegraphics[width=0.95\textwidth]{fig10-mdd-memory.pdf}
\caption{Schematic of a finite-horizon memory-dependent operator and audit questions concerning its kernel, limiting behavior, energy balance, and well-posedness. Model-specific propagation and flux consequences are not asserted; the displayed operator remains subject to primary-source verification.}
\label{fig:mdd_memory}
\end{figure}

\subsection{Notation summary for Section 6}

\begin{table}[htbp]
\centering
\small
\caption{Notation introduced in Section 6.}
\label{tab:notation-sec6}
\begin{tabularx}{\textwidth}{@{}p{1.8cm}X p{2.2cm}p{2.1cm}@{}}
\toprule
\textbf{Symbol} & \textbf{Meaning} & \textbf{Units} & \textbf{Equation} \\
\midrule
$\mathcal{D}[\cdot]$ & Generic memory operator acting on the flux & --- & \eqref{eq:6.1} \\
$\mathcal{D}_{\rm LS}(s)$ & LS transform multiplier (zero initial flux) & dimensionless & \eqref{eq:6.2} \\
${}^C D_t^\alpha$ & Caputo fractional derivative & $t^{-\alpha}$ & \eqref{eq:6.3} \\
${}^{RL} D_t^\alpha$ & Riemann--Liouville derivative & $t^{-\alpha}$ & \eqref{eq:6.4} \\
$a$ & Thermal diffusivity & m$^2$/s & \eqref{eq:6.5} \\
$K(t-\xi)$ & Memory kernel & varies & \eqref{eq:6.6} \\
$D_\omega$ & MDD operator & 1/s & \eqref{eq:6.7} \\
$\omega$ & MDD time delay (horizon) & s & \eqref{eq:6.7} \\
$\kappa(t-\xi)$ & MDD kernel & dimensionless & \eqref{eq:6.8} \\
$\alpha, \beta$ & MDD kernel parameters & varies & \eqref{eq:6.8} \\
$\tilde{q}_i(s)$ & Laplace transform of the heat flux & --- & \eqref{eq:6.2} \\
\bottomrule
\end{tabularx}
\end{table}

\subsection{Open questions for Section 6}

\begin{enumerate}[label=\textbf{OQ6.\arabic*}]
\item What is the complete time-domain energy flux for MDD media with convolution-type memory, expressed solely in terms of the temperature and the stored internal variables? \\ \label{oq:6.1}
\item How is energy partitioned between storage and dissipation in hereditary media, and what is the correct effective flux that enters a local balance? \\ \label{oq:6.2}
\item What are the well-posedness conditions for fractional-order thermoelasticity with arbitrary $\alpha$, in the coupled mechanical setting rather than in isolation? \\ \label{oq:6.3}
\item How do different MDD kernel choices affect the energy flux and the wave propagation, and can the dispersion of the thermal branch be used to identify the kernel experimentally? \\ \label{oq:6.4}
\item Can the memory operator be given a finite-dimensional state-space realisation for a general admissible kernel, and what is the numerical cost of that realisation relative to a direct quadrature? \\ \label{oq:6.5}
\end{enumerate}

% LocalWords:
